\section{Relation to earlier constructions}
\label{sec:finite-prior-comparison}
The determinant $\det[1,v,\ldots,v^{d-1}]$ is a classical local
index form. Arpin, Bozlee, Herr and Smith construct schemes of
monogenerators and polygenerators for finite locally free algebras
\cite{ABHS}. Their Proposition 3.6 and Definition 3.12 identify the
index-form open and its closed complement; Proposition 3.14 gives
the polygenerator open by minors. Thus neither the generator functor
nor its index determinant is claimed as new here. Our use concerns
the coherent cokernel of a global section-multiplication map and,
for generating hyperplanes, its exact closed incidence with a
rank-two quotient. Lemma~\ref{lem:hyperplane-subalgebra} and
Theorem~\ref{thm:universal-hyperplane} prove that closed-scheme
identification on nonreduced bases, rather than deriving it from
agreement of generator opens.

The algebra $G_r$ is the classical two-row Grassmannian cohomology
ring; \cite[Theorem 2.7]{Grinberg} provides the rectangular Schur
basis in a more general quotient. Its role here is an exact
transverse algebra of multiplication failure, determined by a
rank-two quotient of $\C[z]/(z^r)$. The explicit residual equations
\eqref{eq:hyperplane-residual-generators} verify this identification
without a tangent-cone approximation. The diagonal-ideal and
all-powers identities in the three-plane result are precisely
\cite[Corollary 3.8.3]{Haiman}; the argument here supplies the
finite total-degree presentation and conductor transport, not a
new proof of that theorem.

Normal generation in sufficiently positive degree on a curve is
classical \cite{GL86}. The curve conductor theorem uses it only for
$L(-D)$, and then proves the precise quotient comparison by filling
the double-vanishing space and its first normal layer. On $\Pj^1$
the evaluation-kernel argument gives the sharper rank-dependent
bound and its boundary example.

For pencils, confluent Vandermonde factorization and the discriminant
identity are classical algebra. The moving normalization is the
explicit scalar-pair incidence of that index-form divisor. Its
finite map and differential test are proved in
Section~\ref{sec:moving-contact-normalization}; the normalization
is not offered as a construction independent of index-form geometry.
The contact-hyperplane incidence is different: it is the failure
scheme itself, not just its normalization, and its total space is
smooth even at collisions.

Ballico's failure-locus work \cite{Ballico93,Ballico96} is a necessary
nearest-source comparison. The accessible statements of
\cite[Theorem 0.2, Proposition 2.2, Remark 2.3]{Ballico96} concern
failure on finite linear sections of a fixed completely embedded
variety. The complete theorem pages of \cite{Ballico93} have not
been obtained in this revision, so their exact relation to the
present scheme statements remains unverified. We make no inference
of non-anticipation from that access limitation. The results here
are submitted with complete proofs and explicit classical inputs;
priority of the resulting package is a separate question for source
comparison and editorial assessment.
